\clearpage
\fancyhead[L]{\small PIAZZI. STAR CATALOGUE.}
\hypertarget{native-595-lower}{}
\begin{center}\small Göttingische gelehrte Anzeigen. Issue 91. Pages 903..904. 10 June 1815.\\[3pt]\rule{.72\linewidth}{.3pt}\end{center}

\emph{Praecipuarum stellarum inerrantium positiones mediae ineunte saeculo XIX ex observationibus habitis in specula Panormitana ab anno} 1792 \emph{ad annum} 1813. \emph{Ex regia typographia militari.} Palermo 1814. 187 pages in folio. (Price in Florence at Molini, Landi and Co., 30 lire.)

In the years since the appearance of \textsc{Piazzi}’s famous star catalogue of 1803, the fruit of ten years of uninterrupted hard work, observations of the fixed stars at the Palermo Observatory had continued steadily, with \textsc{Piazzi} himself making some and his skilled assistant \textsc{Nicolao Cacciatore} making others. They broadened the observations and even put them on a new footing by establishing the fundamental right ascensions of $\alpha$ in Aquila and $\alpha$ in Canis Minor by direct comparison with the Sun. The new catalogue now presents the full results of all twenty years of work, with the positions of 7646 fixed stars determined to the full precision that observational astronomy can attain today.

In this classic work, \textsc{Piazzi} also devoted particular care to the proper motion of the stars. He relied chiefly on the observations of \textsc{Mayer} and \textsc{Bradley}, and very often also drew on comparisons of his own obser\SourceLeaf{596}vations with one another. The positions of 1041 stars for 1756 are appended at the end of the work, as another of \textsc{Piazzi}’s assistants, \textsc{Joseph Pilati}, reduced them for this purpose from \textsc{Bradley}’s observations.

The annual precession of the stars was calculated using the following formulas.
\begin{center}\small\renewcommand{\arraystretch}{1.18}\begin{tabular}{@{}l@{\hspace{8pt}}l@{}}
in right ascension $.\;.\;.\;.\;.\;.$ & $46''0395+20''0642\sin\alpha\operatorname{tang}\delta$ \\
in declination $.\;.\;.\;.\;.\;.\;.\;.\;.\;.$ & $20''0642\cos\alpha$ \\
\end{tabular}\end{center}
\noindent These numbers are based on the annual lunisolar precession in longitude $50''388$, the annual motion of the equinoctial points on the equator due to planetary effects $0''1814$, and the mean obliquity of the ecliptic for 1800, $23^\circ\;27'\;55''5$. These determinations agree very closely with those of \textsc{Bessel}. (See issue 19 of this journal for 1814 [p. 561 of this volume]).

The outward appearance of the work is certainly worthy of its contents, although it has been printed somewhat more economically than the catalogue of 1803. The right ascensions are given precisely only in arc. In time they are given only to the minute. The stars are all ordered by right ascension and divided into 24 hours, with consecutive numbers in each hour. For every star that has a name of its own, the name is given, chiefly from \textsc{Ulugh Beigh}’s catalogue, while all other designations are omitted for those stars. We do not like this last choice, and we believe all astronomers will find it inconvenient to encounter only the Arabic names, almost all unfamiliar to them, in place of the familiar \textsc{Bayer} letters and \textsc{Flamstead} numbers.

At the end of each hour there are always extensive notes on many of the stars listed in it.

After the catalogue had been fully completed and even printed, \textsc{Piazzi} went on to compare it with the catalogues of zodiacal stars by \textsc{Zach} and \textsc{Barry}. Agreement with the former was very close almost throughout. The Mannheim declinations, by contrast, showed more frequent and larger differences. Only the largest of these, amounting to whole minutes of arc, are given here.

